% !TeX root = ../navier-stokes-manuscript.tex
% ============================================================
% 2.1 Vortex Stretching
% ============================================================

Axial stretching lengthens one narrow material vortex-tube segment along its axis and, in the locally axisymmetric aligned case, contracts its cross-section, draws its rotating fluid inward, and makes the interior spin faster; this spin-up is the positive stretching contribution, while the full signed balance with viscosity decides whether its rotation strengthens overall.
The velocity field carrying this segment is governed by the momentum equation
\[
  \underbrace{\partial_t u}_{\text{local acceleration}}
  +
  \underbrace{(u\cdot\nabla)u}_{\text{nonlinear transport}}
  =
  \underbrace{-\nabla p}_{\text{pressure}}
  +
  \underbrace{\nu\Delta u}_{\text{viscous diffusion}}.
\]
Taking its curl exposes vorticity stretching within nonlinear transport, where it can be measured against viscous diffusion for the same moving segment.

\subsection{Vortex Stretching}\label{sub:ThePerfectlyStretchingVortex}

Let $e$ denote the segment's axis and $L(t)$ its length.
When the strain is locally uniform across the segment and aligned with that axis, its axial extension is
\[
  S:=\frac12\bigl(\nabla u+(\nabla u)^{\mathsf T}\bigr),
  \qquad
  Se=\sigma(t)e,
  \qquad
  \sigma(t)>0,
  \qquad
  \frac{\dot L(t)}{L(t)}
  =
  e\cdot Se
  =
  \sigma(t).
\]
The same extension narrows the segment because incompressibility fixes its material volume $\mathcal V$.
With $A(t):=\mathcal V/L(t)$ as its effective transverse area and $r_{\mathrm{eff}}(t):=\sqrt{A(t)/\pi}$ as its effective radius, the contraction satisfies
\[
  \nabla\cdot u=0,
  \qquad
  \frac{d}{dt}(A L)=0,
  \qquad
  \frac{\dot A}{A}=-\sigma(t),
  \qquad
  \frac{\dot r_{\mathrm{eff}}}{r_{\mathrm{eff}}}
  =
  -\frac{\sigma(t)}{2}.
\]
Thus its rotating fluid is drawn toward the axis as its length increases.
For $\omega:=\nabla\times u$, the curl of the momentum equation gives the vorticity balance carried by the segment,
\[
  \frac{D\omega}{Dt}
  =
  S\omega+\nu\Delta\omega,
  \qquad
  \frac{D}{Dt}
  :=
  \partial_t+u\cdot\nabla,
\]
and the alignment $\omega=|\omega|e$ turns that balance into
\[
  S\omega
  =
  \sigma(t)\omega,
  \qquad
  \frac12\frac{D|\omega|^2}{Dt}
  =
  \sigma(t)|\omega|^2
  +
  \nu\,\omega\cdot\Delta\omega.
\]
The positive term quantifies the spin-up caused by axial stretching, but the signed viscous contribution still decides whether the rotation grows.
The energy identity proved in \Cref{prop:ClassicalKineticEnergyIdentity}, together with the antisymmetric part of $\nabla u$, gives
\[
  \norm{u(t)}_2
  \leq
  \norm{u_0}_2
  <\infty,
  \qquad
  \left|\frac12\bigl(\nabla u-(\nabla u)^{\mathsf T}\bigr)\right|_{\mathrm F}
  =
  \frac{|\omega|}{\sqrt2}
  \leq
  |\nabla u|_{\mathrm F}.
\]
Finite kinetic energy is compatible with local concentration of the rotational gradient in the narrowing segment, so it does not settle the competition at finer scale.

\divider{\NavierStokes{} Scaling}

Fix $t_0<T_*$ and write $\ell:=r_{\mathrm{eff}}(t_0)$ for the segment's effective radius at that time.
The exact full-field comparison at the finer width $\lambda^{-1}\ell$, for $\lambda>1$, is
\[
  u_\lambda(x,t)
  :=
  \lambda u(\lambda x,\lambda^2t),
  \qquad
  p_\lambda(x,t)
  :=
  \lambda^2p(\lambda x,\lambda^2t).
\]
At time $\lambda^{-2}t_0$, the corresponding segment in $u_\lambda$ has effective radius $\lambda^{-1}\ell$.
This rescaled field is another solution at another scale, not a later material state of the opening segment.
Rescaling acts directly on every momentum term:
\[
\begin{aligned}
  \partial_tu_\lambda(x,t)
  &=
  \lambda^3(\partial_tu)(\lambda x,\lambda^2t),
  \\
  (u_\lambda\cdot\nabla)u_\lambda(x,t)
  &=
  \lambda^3\bigl((u\cdot\nabla)u\bigr)(\lambda x,\lambda^2t),
  \\
  \nabla p_\lambda(x,t)
  &=
  \lambda^3(\nabla p)(\lambda x,\lambda^2t),
  \\
  \nu\Delta u_\lambda(x,t)
  &=
  \lambda^3\nu(\Delta u)(\lambda x,\lambda^2t).
\end{aligned}
\]
Viscosity and nonlinear transport acquire the same cubic factor, as do local acceleration and pressure, so finer scale gives viscosity no relative advantage.

\divider{Critical Height}

The unchanged equation balance leaves the homogeneous Sobolev order to distinguish the sizes of the rescaled fields.
Let $\Lambda:=(-\Delta)^{1/2}$.
The Fourier transform obeys
\[
  \widehat{u_\lambda}(\xi,t)
  =
  \lambda^{-2}\widehat u(\lambda^{-1}\xi,\lambda^2t),
\]
and a change of variables then yields
\[
  \norm{u_\lambda(t)}_{\dot H^s}^2
  =
  \lambda^{2s-1}
  \norm{u(\lambda^2t)}_{\dot H^s}^2.
\]
The exponent $2s-1$ is negative at Sobolev order zero, positive at order one, and vanishes at order $s=\tfrac12$.
Define the critical height by
\[
  H(t)
  :=
  \frac12\norm{u(t)}_{\dot H^{1/2}}^2
  =
  \frac12\norm{\Lambda^{1/2}u(t)}_2^2.
\]
For the same quantity $H_\lambda$ evaluated on $u_\lambda$, it and its neighboring norms transform as
\[
  \norm{u_\lambda(t)}_2^2
  =
  \lambda^{-1}\norm{u(\lambda^2t)}_2^2,
  \qquad
  H_\lambda(t)=H(\lambda^2t),
  \qquad
  \norm{\nabla u_\lambda(t)}_2^2
  =
  \lambda\norm{\nabla u(\lambda^2t)}_2^2.
\]
The rescaled solution has less kinetic energy and a larger first-derivative norm, while its half-derivative height is scale neutral.
Its time sign must still come from the signed balance of nonlinear production and viscous dissipation.

\divider{Nonlinear Production versus Viscous Dissipation}

At critical height, the nonlinear contribution is the signed production
\[
  \mathcal P_{1/2}[u](t)
  :=
  -\operatorname{Re}
  \pair
    {\Lambda^{1/2}u(t)}
    {\Lambda^{1/2}\bigl((u(t)\cdot\nabla)u(t)\bigr)}_{L^2}.
\]
Incompressibility removes the pressure pairing, while the Laplacian supplies the positive dissipation on the left of
\[
  H'(t)
  +
  \nu\norm{\Lambda^{3/2}u(t)}_2^2
  =
  \mathcal P_{1/2}[u](t).
\]
The critical height rises exactly when signed nonlinear production exceeds viscous dissipation.
Rescaling multiplies both quantities according to
\[
  \mathcal P_{1/2}[u_\lambda](t)
  =
  \lambda^2\mathcal P_{1/2}[u](\lambda^2t),
  \qquad
  \nu\norm{\Lambda^{3/2}u_\lambda(t)}_2^2
  =
  \lambda^2\nu\norm{\Lambda^{3/2}u(\lambda^2t)}_2^2.
\]
Neither gains ground because they share the same quadratic factor, and this scale match alone does not determine the physical duration of a transition between widths.

\divider{The Heat Clock}

The heat equation $v_t=\nu\Delta v$ supplies a linear clock for that duration through its Fourier modes:
\[
  \widehat v(\xi,t+\delta t)
  =
  e^{-\nu|\xi|^2\delta t}\widehat v(\xi,t).
\]
A scale-localized velocity structure of width $r$ has modes concentrated near $|\xi|\simeq r^{-1}$, and those modes undergo an order-one heat change when
\[
  \nu|\xi|^2\delta t\simeq1,
\]
so its viscous comparison time is
\[
  \tau_\nu(r)
  :=
  \frac{r^2}{\nu}.
\]
Halving the structure's width quarters this clock, and at general scale
\[
  \tau_\nu(\lambda^{-1}r)
  =
  \lambda^{-2}\tau_\nu(r).
\]
The heat clock thus carries the same time contraction as the full \NavierStokes{} rescaling without asserting that one nonlinear field realizes the transition.

\divider{The Zeno Cascade}

Consider the dyadic widths and their heat times
\[
  r_n:=2^{-n}r_0,
  \qquad
  \tau_n:=\frac{r_n^2}{\nu},
\]
which form the finite geometric sum
\[
  \tau_n
  =
  \frac{r_0^2}{\nu}4^{-n},
  \qquad
  \sum_{n=0}^{\infty}\tau_n
  =
  \frac{4r_0^2}{3\nu}
  <
  \infty.
\]
This finite sum provides a possible finite-time cascade schedule only if one \NavierStokes{} velocity field exhibits scale-localized structures at widths
\[
  r_0>r_1>r_2>\cdots\longrightarrow0
\]
and at sampled times
\[
  t_0<t_1<t_2<\cdots\uparrow T_*<\infty,
  \qquad
  t_{n+1}-t_n\asymp\tau_n,
\]
with comparison constants independent of $n$.
Under that condition, the remaining time satisfies
\[
  T_*-t_n
  \asymp
  \sum_{k=n}^{\infty}\tau_k
  =
  \frac{4r_n^2}{3\nu}.
\]
Both the next sampled gap and the total time left are then of order $r_n^2/\nu$.
Even as these intervals collapse, the finite schedule does not yet control the first or any successive Eulerian time derivative of that same field at one fixed spatial point.
